% Generated draft: requires mathematical and source review.
\documentclass[11pt]{article}
\usepackage[margin=25mm]{geometry}
\usepackage{fontspec}
\setmainfont{Latin Modern Roman}
\usepackage{amsmath,amssymb,mathtools,unicode-math}
\setmathfont{Latin Modern Math}
\usepackage{enumitem,needspace,xurl,hyperref}
\setlength{\parindent}{0pt}
\setlength{\parskip}{5pt}
\setlength{\emergencystretch}{3em}
\newcommand{\N}{\mathbb{N}}
\newcommand{\Z}{\mathbb{Z}}
\newcommand{\Q}{\mathbb{Q}}
\newcommand{\R}{\mathbb{R}}
\newcommand{\C}{\mathbb{C}}
\newcommand{\F}{\mathbb{F}}
\newcommand{\A}{\mathbb{A}}
\newcommand{\PP}{\mathbb{P}}
\newcommand{\E}{\mathbb{E}}
\newcommand{\eps}{\varepsilon}
\newcommand{\dd}{\,\mathrm{d}}
\begin{document}
\section*{UnsolvedMath ID: 30000624}
\textbf{Source problem code:} catalogue entry
\subsection{Definitions and mathematical statement}
Let $L$ be a centreless Lie torus in the sense intended by the source. Centreless means
\[
 Z(L)=\{x\in L:[x,L]=0\}=0.
\]
The Lie-torus axiom system, base field, root system, external grading group, and finiteness assumptions are not supplied. In the usual framework a Lie torus has compatible root-lattice and external-group gradings with specified root-space and division properties; no particular version of those axioms is silently selected.

For an extended affine Lie algebra $E$ with its distinguished Cartan and invariant-form data, its core $E_c$ is conventionally the ideal generated by the nonisotropic root spaces. Its centreless core is
\[
 E_{cc}=E_c/Z(E_c).
\]
The source asks whether
\[
 \forall L\text{ centreless Lie torus}\ \exists E\text{ extended affine Lie algebra}:
       E_c/Z(E_c)\cong L,
\]
and, if so, how $E$ can be reconstructed from $L$. The quotient is by the center of the core, not by substituting $E/Z(E)$.

A precise reconstruction must specify the additional central-extension, derivation, invariant-form, or other data required by the source's framework and the equivalence under which reconstructions are compared. The question does not assert that $L$ determines a unique $E$, nor that every compatible choice gives one. It also does not state whether the displayed isomorphism preserves the Lie-torus gradings or is only an abstract Lie-algebra isomorphism.

The exact extended-affine axioms, nullity or rank restrictions, and base-field assumptions must be recovered. No finite-generation-over-the-centroid condition, characteristic-zero hypothesis, or narrower torus class is imposed as an unstated source fact. This retains both universal existence and the constructive/reconstruction task without inventing a classification theorem.
\subsection{Short English statement}
The question asks whether every centreless Lie torus is the core modulo its own center of an extended affine Lie algebra, and what data construct such an algebra from the torus. The precise Lie-torus/EALA classes and graded equivalence are unspecified; uniqueness is not asserted.
\subsection{Sources}
\begin{itemize}
\item Local UnsolvedMath registry, permanent ID 30000624, code catalogue entry, ``Reconstruction of Extended Affine Lie Algebras from Lie Tori.'' Proposed year 2006; proposer unspecified.
\item \emph{Infinite-Dimensional Lie Theory} (2007), Oberwolfach Reports, pp.\ 3259--3320, DOI 10.4171/owr/2006/55; \url{https://doi.org/10.4171/owr/2006/55}. The displayed report year 2007 differs from the DOI and proposed year 2006. This is the sole supplied literature reference.
\item Statement status corrected\_verified; the supplied note says editorial and LaTeX cleanup preserved meaning without added source-dependent mathematical content.
\item Literature triage checked 2026-08-21: open. It generically calls the item an unresolved problem/conjecture/direction and reports no authoritative complete resolution in the supplied targeted check. Exact axiom conventions and current status require independent verification. Created 2026-08-25; updated 2026-09-16.
\item Registry curation by Ulam AI / UnsolvedMath Contributors, CC BY 4.0; the underlying report retains its supplied provenance. Unreviewed statement-only draft.
\end{itemize}
\end{document}
